% introduction.tex
% EXEMPLAR section showing the patterns used throughout the template:
%   - top comment header naming the section and (where applicable) the
%     experiments / derivations it depends on
%   - \label{sec:...} immediately after the section start in main.tex
%   - subsection labels of the form \label{subsec:...}
%   - cross-references via \ref{} / \autoref{}
%   - figure include via \input{...} of a figure-fragment file
%   - bibliography citations with \cite{key}
%   - audit-table reference to Table~\ref{tab:audit}
%
% Replace the prose with your own; keep the labelling and cross-reference
% conventions.

\label{sec:introduction}

\placeholder{One- or two-paragraph orientation: what problem the paper
addresses, why existing approaches fall short, and what the paper's
contribution is. The introduction should be readable without the reader
having internalised any framework-specific notation yet.}

\subsection{Scope}
\label{subsec:scope}

\placeholder{What is in scope for this paper, and what is explicitly
deferred to companion papers or follow-on work. Reference the audit
table (Table~\ref{tab:audit}) if you maintain one --- it is the
canonical answer to ``has this claim been verified yet?''}

%% The audit table is included here so the reader meets it before any
%% claim that depends on a status entry.  Move it elsewhere (e.g. into
%% a "Predictions" section) if that fits the paper better.  Delete this
%% \input if you are not maintaining an audit table.
% ============================================================
% audit_table.tex
%
% A "system audit" table mapping the paper's claims to the evidence
% backing each one.  This is *the* canonical source of truth for which
% claims are PASS / PART / STUB / FAIL --- prose elsewhere in the paper
% should not contradict it.
%
% Status semantics (paste into the abstract or front-matter as well):
%   PASS  -- derivation complete or experiment confirms the claim to
%            stated precision.
%   PART  -- mechanism demonstrated, full quantitative match pending;
%            specific gaps documented in the evidence column.
%   STUB  -- analytical result with no numerical confirmation yet,
%            or planned experiment not yet run.
%   FAIL  -- tested and disconfirmed.
%
% Use \texttt{} for code/experiment names and file paths inside cells.
% Long cells are fine -- the column widths are tuned for prose.
% ============================================================

\label{tab:audit}

{\scriptsize
\setlength{\tabcolsep}{4pt}
\renewcommand{\arraystretch}{0.95}
\begin{longtable}{@{}p{2.6cm}p{3.8cm}p{3.2cm}p{4.4cm}p{1.0cm}@{}}
\caption{\small System audit: paper claims mapped to supporting
derivations and experiments.  Status semantics in the front-matter
(\texttt{PASS} / \texttt{PART} / \texttt{STUB} / \texttt{FAIL}).}
\label{tab:audit_caption} \\
\hline
\textbf{Claim} & \textbf{Mechanism} & \textbf{Continuum / formal limit} & \textbf{Evidence} & \textbf{Status} \\
\hline
\endfirsthead

\multicolumn{5}{l}{\textit{(Table~\ref{tab:audit}, continued from previous page)}} \\
\hline
\textbf{Claim} & \textbf{Mechanism} & \textbf{Continuum / formal limit} & \textbf{Evidence} & \textbf{Status} \\
\hline
\endhead

\hline
\multicolumn{5}{r}{\textit{(continued on next page)}} \\
\endfoot

\hline
\endlastfoot

% ── Example row 1 (replace with a real claim) ───────────────────────
Example claim, fully derived
    & One-line summary of the mechanism
    & The standard / continuum statement it reduces to
    & Derived analytically (\S\ref{subsec:cross_ref_example}, Eq.~\ref{eq:example_identity})
    & \texttt{PASS} \\

% ── Example row 2 (replace with a real claim) ───────────────────────
Example claim, partial evidence
    & One-line summary
    & The target statement
    & \texttt{exp\_XX}: describe the gap concretely (e.g.\ ``mechanism
      reproduced; quantitative match pending higher-resolution run'')
    & \texttt{PART} \\

\hline
\end{longtable}
}


\subsection{A Worked Cross-Reference Example}
\label{subsec:cross_ref_example}

A statement supported by an external citation looks like this~\cite{example2024}.
A statement supported by an in-paper derivation looks like this
(see~\autoref{subsec:scope}).  An equation that is referenced
elsewhere is given a label:
%
\begin{equation}
    \label{eq:example_identity}
    e^{i\pi} + 1 = 0.
\end{equation}
%
Equation~\eqref{eq:example_identity} is referenced by label, not by
hard-coded number.

\subsection{A Worked Figure-Include Example}
\label{subsec:figure_example}

Figures live in \texttt{paper/figures/} (or in the repo-root
\texttt{figures/} directory if the figure is shared with experiments).
The standard pattern is to keep the \verb|\includegraphics| line and the
caption in a small fragment file under \texttt{paper/figures/}, then
\verb|\input| that fragment from inside a \verb|figure| environment
here.  See \texttt{paper/figures/example\_figure.tex} for the template;
once you replace the placeholder with a real figure, uncomment the
include block below.

% \begin{figure}[H]
%     \centering
%     % example_figure.tex
% Figure-fragment file. Convention: \includegraphics on line 1, then the
% caption and label.  Included from a section via:
%   \begin{figure}[H]
%       \centering
%       % example_figure.tex
% Figure-fragment file. Convention: \includegraphics on line 1, then the
% caption and label.  Included from a section via:
%   \begin{figure}[H]
%       \centering
%       % example_figure.tex
% Figure-fragment file. Convention: \includegraphics on line 1, then the
% caption and label.  Included from a section via:
%   \begin{figure}[H]
%       \centering
%       \input{figures/example_figure}
%   \end{figure}
% Putting the caption next to the includegraphics keeps the figure file
% self-contained -- if you reuse the same plot in a talk, you take this
% one .tex file with you.

\includegraphics[width=0.7\linewidth]{figures/example_figure.pdf}
\caption[Short list-of-figures caption]{%
    Long caption with full prose explanation of axes, units, and what
    the reader should take away from the figure.  This is the version
    that appears under the figure in the body.}
\label{fig:example}

%   \end{figure}
% Putting the caption next to the includegraphics keeps the figure file
% self-contained -- if you reuse the same plot in a talk, you take this
% one .tex file with you.

\includegraphics[width=0.7\linewidth]{figures/example_figure.pdf}
\caption[Short list-of-figures caption]{%
    Long caption with full prose explanation of axes, units, and what
    the reader should take away from the figure.  This is the version
    that appears under the figure in the body.}
\label{fig:example}

%   \end{figure}
% Putting the caption next to the includegraphics keeps the figure file
% self-contained -- if you reuse the same plot in a talk, you take this
% one .tex file with you.

\includegraphics[width=0.7\linewidth]{figures/example_figure.pdf}
\caption[Short list-of-figures caption]{%
    Long caption with full prose explanation of axes, units, and what
    the reader should take away from the figure.  This is the version
    that appears under the figure in the body.}
\label{fig:example}

% \end{figure}
